\documentclass[11pt]{article}
\usepackage{amsmath,amssymb,amsthm,hyperref}
\newtheorem{proposition}{Proposition}
\newtheorem{corollary}{Corollary}
\title{Why irrational two-block quotient modes are simpler}
\author{Sparse partitioning research session}
\date{November 15, 2026}
\begin{document}
\maketitle
\paragraph{Status.}
These are elementary consequences of Perron--Frobenius and Galois
conjugacy, not a claimed new algebraic mechanism or a historical-firstness
claim. The observation was identified during the independent audit of
the unequal-block finite-field theorem. It supplements that theorem
without changing its frozen source or its exact stated scope.

Let $G$ be a connected simple graph with an equitable partition into two
nonempty vertex classes. Write its integer quotient as
\[
 Q=\begin{pmatrix}a_1&b_1\\b_2&a_2\end{pmatrix},
 \qquad b_1,b_2>0,
\]
where $a_i$ is the number of internal neighbors in class $i$ and $b_i$
the number of neighbors in the other class. Necessarily
$n_1b_1=n_2b_2$. The graph's adjacency matrix is denoted by $A$.

\begin{proposition}[Deterministic irrational-mode simplicity]
If $D=(a_1-a_2)^2+4b_1b_2$ is not a perfect square, both quotient
eigenvalues
\[
 \lambda_\pm=\frac{a_1+a_2\pm\sqrt D}{2}
\]
are simple eigenvalues of $A$. In particular the smaller one recovers
the two classes exactly by the signs of its eigenvector. No parity,
internal-degree, or spectral-outlier assumption is required.
\end{proposition}
\begin{proof}
The quotient is irreducible and nonnegative. Its positive Perron vector
lifts to an entrywise positive adjacency eigenvector at $\lambda_+$.
As $G$ is connected, Perron--Frobenius makes this the simple top
eigenvalue of $A$. The nonsquare discriminant means that $\lambda_+$
and $\lambda_-$ have the same monic irreducible quadratic minimal
polynomial over $\mathbb Q$. In the integer characteristic polynomial
of $A$, two roots conjugate over $\mathbb Q$ have equal multiplicities.
Thus simplicity at $\lambda_+$ implies simplicity at $\lambda_-$.

The smaller root is strictly below both $a_1$ and $a_2$. Consequently
its quotient eigenvector
$(b_1,-(a_1-\lambda_-))$ has opposite nonzero signs. It lifts to the
full graph, and simplicity forces its sign partition to be the unique
one obtained from this eigenvalue.
\end{proof}

This does not imply simplicity in the integer case. For example the
connected cube has quotient eigenvalues $3$ and $1$ under any coordinate
bisection, but the adjacency eigenvalue $1$ has multiplicity three.
The irrational case follows from the simple Perron root precisely
because its conjugacy links the two roots. If the discriminant is a
square, both roots are rational algebraic integers and hence integers,
and this link is absent.

\begin{corollary}[Combining with the finite-field result]
Fix $a_1,a_2\geq3$ and $b_1,b_2\geq1$. Join two independent uniformly
random simple internal regular graphs of degrees $a_1,a_2$ by any simple
$(b_1,b_2)$-biregular graph chosen in advance. Then:
\begin{enumerate}
\item If both block sizes are even and tend to infinity at the required
      fixed ratio, both quotient eigenvalues are simple with high
      probability, whether integer or irrational. The integer case is
      the finite-field theorem in \texttt{unequal-block-integer-mode.tex}.
\item If the quotient discriminant is nonsquare, the same conclusion
      holds along ALL admissible block sizes. Independent internal
      uniform regular graphs of fixed degree at least three are
      connected with high probability [W, Section 2.6]. A positive cross degree then
      connects the two classes, and the proposition applies.
\end{enumerate}
In these cases the smaller mode yields exact two-class recovery. No
claim for odd block sizes in the integer-mode case is supplied here.
\end{corollary}

For two classes with the same total degree $d$, the quotient roots are
$d$ and $d-b_1-b_2$, both integers. Thus the nontrivial equal-total-degree
recovery result needs the finite-field argument and is not covered by
the elementary irrational-mode observation. If the total degrees differ,
vertex degrees themselves distinguish the two classes; the simplicity
proposition remains a spectral fact, but it should not be advertised as
a new hard clustering algorithm for that two-class situation.
\begin{thebibliography}{1}
\bibitem{W} N. C. Wormald, \emph{Models of random regular graphs},
author PDF, model/parity conventions in Section 1 and connectivity
statement and proof in Section 2.6.
\url{https://users.monash.edu/~nwormald/papers/regsurvey.pdf}.
\end{thebibliography}
\end{document}
